\section{非线性椭圆方程}\label{chap:nonlinear-elliptic-equations}

本章从三个角度讨论非线性椭圆方程: 隐函数定理, 变分法和非线性算子理论. 这是本书唯一假设读者熟悉测度论基本结果的一章. 特别地, 假设读者理解以下概念和结果.\begin{itemize}\item 零测集的定义, 以及函数 ``几乎处处\textquotesingle\textquotesingle\ 相等的概念.\item Lebesgue 可测函数的概念, 以及将 $L^p$ 空间定义为几乎处处相等的函数所构成的等价类.\item $L^p$ 空间的 ``测度论\textquotesingle\textquotesingle\ 定义与本书其他部分使用的 ``完备化\textquotesingle\textquotesingle\ 定义的等价性.\item 函数序列几乎处处收敛的概念, 以及各种收敛之间的相互关系. 这包括理解 Fatou 引理和 Lebesgue 控制收敛定理等结果.\end{itemize}

\subsection{扰动结果}\label{sec:nonlinear-elliptic-perturbation-results}

许多关于微分方程的结果表明, 只要考虑足够小的解, 使线性项占优于非线性项, 非线性方程的行为就基本上与其线性化相同. 在 Chapter~\ref{chap:introduction} 中陈述了经典微积分中的隐函数定理, 它对有限维方程组提供了这样的结果. 本节将把隐函数定理推广到 Banach 空间的情形, 再考虑其对椭圆 PDE 的应用.

\subsubsection{Banach 压缩原理与隐函数定理}\label{subsec:banach-contraction-implicit-function}

\PDEIndexMark{idx:banach-contraction-principle-introduction}Banach 压缩原理是求解非线性方程最常用的技巧之一. 其内容为以下定理.

\begin{Theorem}{Banach 压缩}
\label{thm:banach-contraction-fixed-point}
设 $(X,d)$ 为完备度量空间. 假设 $X$ 非空, 且 $T:X\to X$ 为压缩映射, 即存在 $\theta\in[0,1)$, 使对每个 $x,y\in X$ 都有 $d(T(x),T(y))\leq\theta d(x,y)$. 则 $T$ 有唯一不动点.
\end{Theorem}

\begingroup
\renewcommand{\endProof}{\par\nobreak\hfill$\blacksquare$\par\vspace{0.1cm}}
\begin{Proof}
假设有两个不动点 $x$ 和 $y$. 则
\begin{equation}\label{eq:banach-contraction-fixed-point-uniqueness}
d(x,y)=d(T(x),T(y))\leq\theta d(x,y),
\end{equation}
这只有在 $x=y$ 时才可能.

现在任取一点 $x_0\in X$, 递归定义 $x_{n+1}=T(x_n)$. 则有 $d(x_{n+1},x_n)\leq\theta d(x_n,x_{n-1})$, 由此容易得到序列 $x_n$ 是 Cauchy 序列. 设 $x$ 为极限. 则有
\begin{equation}\label{eq:banach-contraction-iteration-limit-fixed-point}
T(x)=T\left(\lim_{n\to\infty}x_n\right)=\lim_{n\to\infty}T(x_n)=\lim_{n\to\infty}x_{n+1}=x.
\end{equation}
证明完成.
\end{Proof}
\endgroup
